\begin{explanationbox}
	\textbf{\textcolor{CExplanation}{一句话直觉}}
	排列数公式给出了从 $n$ 个不同元素中选出 $m$ 个并按一定顺序排列的方法总数，它是分步乘法计数原理的直接结果。

	\medskip
	\textbf{\textcolor{CExplanation}{核心拆解}}
	\begin{description}[style=nextline, leftmargin=2.8em, labelsep=0.5em, itemsep=0.45em]
		\item[\textbf{要点一（分步选位）：}]
			第一个位置有 $n$ 种选法，第二个有 $n-1$ 种，依次类推，第 $m$ 个有 $n-m+1$ 种。
		\item[\textbf{要点二（连乘与阶乘）：}]
			上述连乘 $n(n-1)\cdots(n-m+1)$ 正是 $\frac{n!}{(n-m)!}$。
	\end{description}

	\medskip
	\textbf{\textcolor{CExplanation}{几何本质（有序排布）}}
	排列对应着将 $m$ 个不同元素按顺序放入 $n$ 个位置中的 $m$ 个位置，每个排列是一种有序配置。

	\medskip
	\textbf{\textcolor{CExplanation}{代数意义（阶乘运算）}}
	排列数公式是阶乘的变形，$A_n^m = \frac{n!}{(n-m)!}$ 体现了部分排列与全排列的比例关系。

	\medskip
	\textbf{\textcolor{CExplanation}{考点价值（高频计算）}}
	在高考中，排列数公式是计数原理的基础，常出现在选择题或与组合数、概率结合的题目中。

	\medskip
	\textbf{\textcolor{CExplanation}{顿悟点}}
	\textbf{排列与顺序有关。}只要顺序不同就是不同排列，因此计算时直接使用乘法原理即可。

	\medskip
	\textbf{\textcolor{CExplanation}{使用场景}}
	\begin{itemize}[leftmargin=*, itemsep=0.5em, topsep=0.4em]
		\item \textbf{场景一（排队顺序）：} 从若干人中选部分排成一列。
		\item \textbf{场景二（数字组数）：} 用给定数字组成多位数（数字不重复）。
		\item \textbf{场景三（选号排位）：} 从集合中取元素排成序列。
	\end{itemize}

\end{explanationbox}
